\chapter*{Введение}
\addcontentsline{toc}{chapter}{Введение}

\textbf{Актуальность темы исследования.}

В условиях стремительного роста объёмов визуальных данных, поступающих из
распределённых сетей камер наблюдения и иных сенсоров, базовая задача поиска
целевого объекта --- человека --- становится одним из ключевых ограничителей
эффективности прикладных решений. Информационно‑поисковые системы (ИПС),
работающие с видеопотоками, должны обеспечивать обнаружение, сопоставление и
последующую реидентификацию субъектов в динамичной среде, где изменяются
освещение, ракурс, поза, частота окклюзий и плотность сцены. При этом от таких
систем ожидают оперативности (минимум задержек при принятии решений),
масштабируемости (устойчивая работа при росте числа камер и пользователей) и
надёжности (устойчивость к шумам, пропускам и ошибкам промежуточных этапов).
Иными словами, речь идёт не только о распознавании как таковом, но и о поиске в
визуальном пространстве: по запросу (образцу, изображению или сценарию поиска)
необходимо быстро и корректно выделить все релевантные вхождения одного и того
же субъекта среди большого массива гетерогенных видеоданных.

Традиционные конвейеры обработки в ИПС строятся как последовательность этапов:
детектирование человеческих фигур и лиц в потоке изображений, отслеживание
движений объектов в сцене (трекинг), построение признаковых представлений и,
наконец, идентификация/реидентификация по извлечённым признакам. Такая
архитектура естественна и технологически проверена, однако она несёт типичные
риски <<каскадного эффекта>>: неточность детектирования приводит к деградации
качества эмбеддингов, а далее --- к ошибкам сопоставления и поиска. В реальных
условиях это проявляется особенно заметно: отсутствие фронтального лица,
частичная видимость силуэта, пересечения траекторий, вариативность внешнего
вида, изменение числа людей в кадре и плотности сцены приводят к росту ложных
срабатываний и пропусков. Следовательно, задача повышения качества поиска в ИПС
сводится к согласованной оптимизации всего контура обработки --- от первичного
выделения объектов до стратегии слияния многомодальных признаков при принятии
решения.

В современном информационном обществе задачи, связанные с поиском, распознаванием
и идентификацией объектов, приобретают всё большее значение в связи с активным
ростом объёмов визуальных данных, получаемых с цифровых видеокамер и иных
сенсоров, используемых в системах видеонаблюдения. ИПС, ориентированные на поиск
и идентификацию человека, играют важную роль как в области обеспечения
общественной безопасности, так и в интеллектуальных городских системах
мониторинга и управления доступом, в транспортной инфраструктуре, на режимных
объектах, в системах контроля доступа и в бортовых интеллектуальных системах.

Под \textit{детектированием} в настоящей работе понимается процесс автоматического
выделения областей изображения или видеопотока, соответствующих объектам
заданного класса, в частности --- человеческим лицам и силуэтам. Данный этап
критически важен, так как именно качество выделения объектов во многом
определяет общую точность последующих этапов анализа и реидентификации.
В контексте задач идентификации и реидентификации человека детектирование
является первым и ключевым этапом обработки данных: именно от его качества
зависит корректность работы всех последующих шагов, включая построение
признаковых представлений (эмбеддингов) и процедуру сопоставления. Ошибки на
стадии детектирования приводят к каскадному накоплению неточностей и снижению
итоговой точности системы. Для ИПС, работающих в условиях большого потока
данных и разнообразия внешних факторов (освещение, ракурс, плотность сцены,
частичные перекрытия объектов), детектирование особенно важно, так как оно
задаёт верхнюю границу эффективности алгоритмов распознавания.

Существенный пласт исследований в области компьютерного зрения посвящён
обнаружению и распознаванию лиц. Методы обнаружения лиц~\cite{Shanmugavadivu2016facedetect,Hofer2021faceparts,He2020crossspectral,Li2006nirface,Zhang2004infofusion,CLiu2014face} и
методы распознавания лиц~\cite{Zhao2011facecorner,Dang2017facedetection,Tian2012faceJPEG,Nehru2017violajones,CLiu2014face} активно развиваются в течение более
чем двух десятилетий. В современных системах изображения лиц по‑прежнему
остаются предпочтительным источником биометрической информации для
идентификации личности. Вместе с тем изображения лиц доступны не всегда, что
делает автономное распознавание по лицу неприменимым в ряде эксплуатационных
сценариев. Это особенно заметно в уличных и транспортных сценах, при
неблагоприятном ракурсе, неполной видимости, окклюзиях или недостаточном
разрешении лица.

Наряду с лицевой биометрией развивается распознавание человека по нелицевым
признакам. Распознавание человека может быть направлено на идентификацию
личности по изображению или видео без обязательного присутствия лица.
Существует ряд подходов, основанных на анализе отличных от лица признаков
человека. Как показано в работах~\cite{Cutler2000periodicmotion,Heisele1998pedestrian}, периодичность походки позволяет
обнаруживать идущего человека в последовательности изображений. Походка также
использовалась для распознавания людей в видеопоследовательностях. Однако, как
отмечалось в~\cite{Nixon1999gait}, подобные подходы часто тестировались на ограниченном числе
людей, и их переносимость на большие совокупности субъектов остаётся
неоднозначной.

На фоне ограничений автономного распознавания по лицу или по силуэтным
признакам всё большую роль играют методы \textit{реидентификации человека}
(ReID). Реидентификация человека представляет собой процесс определения и
подтверждения личности индивида на основе сравнения его характеристик,
полученных из различных источников данных или в разные моменты времени.
В контексте компьютерного зрения и обработки изображений речь идёт об
идентификации одного и того же человека на разных изображениях или видеозаписях
при изменении условий съёмки (освещение, ракурс, поза, фон, одежда и др.).
Последние работы по person search и person re-identification условно делятся
на одноэтапные методы~\cite{Chang2018rcaa,Chen2020hierarchical,Chen2020normaware} и двухэтапные методы~\cite{Chen2020personsearch,Dong2020personsearch,Han2019reidlocalization}. Типичный
двухэтапный метод сначала обнаруживает людей, а затем применяет модель
распознавания к обнаруженным изображениям. Напротив, одноэтапные методы
оптимизируют задачи обнаружения и распознавания людей одновременно. Чтобы
повысить различительную способность моделей, в некоторых работах вводятся
вспомогательные задачи~\cite{Chen2020personsearch,Chen2022keypoint} (например, оценка позы, распознавание атрибутов,
семантическая сегментация). В~\cite{Zhong2020personsearch} извлекаются признаки видимых частей тела и
выполняется сопоставление частичных признаков. В~\cite{Chen2020personsearch} исследуется влияние
справочной информации и структур сцены на поиск человека. В рамках отечественных
исследований проблема реидентификации человека активно разрабатывается рядом
научных школ, включая работы Зайцева А.\,Ю., Сахарова В.\,Н., Лапшина С.\,А.
и других, в которых рассматриваются как алгоритмические, так и архитектурные
аспекты построения систем идентификации в сложных условиях наблюдения~\cite{Zaitsev2021reid,Sakharov2020recognition,Lapshin2019identification,Smirnov2022gait,Gorelov2020review}.

С практической точки зрения ИПС должны не только находить человека <<здесь и
сейчас>>, но и сопоставлять наблюдения во времени и пространстве, обеспечивая
корректное <<сведение>> следов субъекта под изменяющимися условиями. Отсюда
вытекают прикладные требования к таким системам:
\begin{itemize}
    \item инвариантность к внешним факторам (освещение, ракурс, поза, частичная
    окклюзия, изменчивость одежды и плотности сцены);
    \item контролируемая вычислительная сложность для работы в режиме, близком
    к реальному времени;
    \item формализуемые метрики качества (mAP, CMC/Rank, ROC/PR) с привязкой
    к эксплуатационным сценариям (порог решения, режим Top-k, режим воздержания
    при низкой уверенности);
    \item технологическая интегрируемость в уже развёрнутые платформы наблюдения
    и аналитики без разрушения существующих регламентов хранения и обработки данных.
\end{itemize}

Отдельного внимания заслуживает текущая практика внедрения биометрических
систем и коммерческих решений. Рост рынка биометрических технологий и широкое
распространение систем распознавания лица в России и мире~\cite{Nazarova2023biometry,Mubinova2024face,Emets2019biometry}
показывают высокую востребованность соответствующих методов. Вместе с тем
использование коммерческих решений нередко осложняется высокой стоимостью,
закрытостью алгоритмов, непрозрачностью метрик качества и отсутствием
возможности тонкой адаптации под конкретные сцены и условия эксплуатации.
Таким образом, сохраняется запрос на научно обоснованные и воспроизводимые
методы, которые могли бы быть интегрированы в реальные ИПС и поддерживали бы
контролируемую точность, устойчивость и интерпретируемость принятия решений.

Несмотря на большое количество подходов, существующие ИПС распознавания и
реидентификации человека обладают рядом ограничений. Во‑первых, большинство
существующих решений базируется на использовании одного типа признаков
(лицевые либо нелицевые), что не позволяет учитывать многофакторную природу
идентификации личности. Такая однофакторность приводит к снижению точности,
особенно при отсутствии лица или его частичном перекрытии. Во‑вторых, многие
подходы не обеспечивают эффективной работы в условиях динамично изменяющихся
сцен, где количество людей и условия видимости варьируются во времени.
В‑третьих, существенной проблемой остаётся неустойчивость алгоритмов в условиях
неопределённости: изменения освещения, разнообразия поз, динамики передвижений,
появления и исчезновения модальностей, взаимодействия множества объектов
в сцене.

Поэтому актуальной является разработка комплексной методики реидентификации
человека, учитывающей многофакторность процесса и способной адаптироваться
к изменяющимся условиям наблюдения. В настоящей работе используется подход,
в котором ключевую роль играет согласованное использование двух основных
модальностей --- лицевой и силуэтной. Такой подход опирается не на механическое
усреднение представлений, а на формализованное объединение источников признаков
с учётом их текущей доступности и надёжности.

В общем случае состояние наблюдаемого субъекта может быть описано вектором
модальностей
\begin{equation}
    \mathbf{z}(t)=\bigl(\mathbf{z}_{f}(t),\,\mathbf{z}_{s}(t),\,\mathbf{z}_{a}(t),\,\ldots\bigr),
\end{equation}
где $\mathbf{z}_{f}$ --- лицевые признаки, $\mathbf{z}_{s}$ --- признаки
силуэта, $\mathbf{z}_{a}$ --- дополнительные атрибутивные признаки, а
многоточие отражает возможность включения других биометрических и контекстных
факторов. В существующих системах, как правило, используется только одна
компонента, чаще всего лицевые признаки, что фактически редуцирует
многофакторную задачу к однофакторной постановке и делает систему уязвимой
к отсутствию соответствующей модальности.

В рассматриваемых эксплуатационных сценариях доступный набор модальностей
изменяется во времени:
\begin{equation}
    \mathcal{M}(t)\subseteq\{f,s,a,\ldots\},
\end{equation}
и в отдельные моменты времени некоторая компонента, например $f$, может
отсутствовать вследствие окклюзий, поворота головы или недостаточного
разрешения. Следовательно, методика реидентификации должна не только объединять
несколько групп признаков, но и адаптивно учитывать их текущую доступность
и надёжность, перераспределяя вклад модальностей в итоговое решение в
зависимости от условий наблюдения.

Актуальность разработки именно многофакторной и адаптивной методики
обусловлена тем, что ни один из факторов по отдельности не обеспечивает
требуемого уровня точности в реальных условиях эксплуатации. При фиксированном
использовании только лицевых признаков качество резко деградирует при частичном
или полном отсутствии лица, тогда как использование только силуэта и атрибутов
недостаточно для разграничения большой совокупности людей с близким внешним
видом. Совместный учёт нескольких комплементарных факторов позволяет, во‑первых,
компенсировать отсутствие или низкое качество отдельной модальности,
а во‑вторых, формировать более информативное признаковое представление,
непосредственно приводящее к снижению вероятности ошибки реидентификации.

Таким образом, актуальность настоящего исследования определяется необходимостью
повышения эффективности процессов распознавания и реидентификации людей
в информационно‑поисковых системах в условиях информационной неопределённости
и изменяющихся факторов наблюдения при рациональных вычислительных затратах.

\textbf{Степень разработанности темы.}

Задачи распознавания человека по изображению и видео, включая детектирование,
идентификацию и реидентификацию (ReID), относятся к числу наиболее активно
развивающихся направлений цифровой обработки аудиовизуальной информации и
распознавания образов. Для распознавания лица исторически сформировались
фундаментальные подходы выделения признаков и снижения размерности, включая
PCA/<<eigenfaces>> и LDA, а также различные схемы каскадного детектирования,
локализации лицевых ключевых точек и последующего сопоставления признаков.
В более поздних работах широкое распространение получили глубокие нейронные
сети, обеспечившие качественный переход от ручных признаков к обучаемым
эмбеддингам~\cite{Shanmugavadivu2016facedetect,Hofer2021faceparts,He2020crossspectral,Li2006nirface,Zhang2004infofusion,CLiu2014face,Zhao2011facecorner,Dang2017facedetection,Tian2012faceJPEG,Nehru2017violajones}. Существенный прогресс
в развитии дискриминативных эмбеддингов связан с метрическим обучением
и функциями потерь, формирующими выраженные межклассовые границы
в признаковом пространстве, в том числе ArcFace.

Распознавание человека по нелицевым признакам также имеет долгую историю.
Помимо походки~\cite{Cutler2000periodicmotion,Heisele1998pedestrian,Nixon1999gait}, используются силуэт, атрибуты внешнего вида,
семантическая сегментация частей тела, контекст и структура сцены. Эти
направления важны, поскольку лицо не всегда доступно, а для массовых сцен,
нефронтальных ракурсов и больших расстояний до камеры нелицевые признаки
становятся основным источником информации о личности.

Современное направление person re-identification сформировалось как ответ
на практические задачи мультикамерного наблюдения и поиска людей. Обзоры
и систематизации области~\cite{BedagkarGala2014survey,Wu2019reidsurvey,Zheng2016reidsurvey,Vezzani2013reidsurvey,Wang2020reidoverview,Yadav2023reidsurvey,Ye2021reidsurvey} показывают, что сегодня
ReID включает широкий спектр постановок и подходов: двухэтапные и одноэтапные
схемы, методы person search, part-based and attention-based модели,
многомодальные подходы, методы partial matching, re-ranking, графовые и
структурные модели. При этом в реальных ИПС качество итогового результата
по‑прежнему чувствительно к ошибкам ранних стадий обработки и к вариативности
условий наблюдения.

В работах~\cite{Chang2018rcaa,Chen2020hierarchical,Chen2020normaware} рассматриваются одноэтапные методы поиска человека,
одновременно решающие задачи обнаружения и распознавания. В~\cite{Chen2020personsearch,Dong2020personsearch,Han2019reidlocalization}
представлены двухэтапные схемы, в которых сначала выполняется локализация
человека, а затем --- извлечение признаков и сопоставление. В ряде работ
используются вспомогательные задачи~\cite{Chen2020personsearch,Chen2022keypoint}, например оценка позы и распознавание
атрибутов, а также механизмы выделения видимых частей тела~\cite{Zhong2020personsearch}. Всё это
подтверждает, что область ReID движется в сторону более сложных и информативных
многокомпонентных моделей.

В отечественной научной повестке задача реидентификации и построения систем
идентификации в сложных условиях наблюдения разрабатывается рядом научных школ
и исследовательских коллективов. В работах~\cite{Zaitsev2021reid,Sakharov2020recognition,Lapshin2019identification,Smirnov2022gait,Gorelov2020review} рассматриваются как
алгоритмические, так и архитектурные аспекты построения систем идентификации
человека по видеоизображениям, а также комбинированные решения, учитывающие
лицевые, силуэтные и иные биометрические признаки. Эти исследования подтверждают
актуальность темы и наличие устойчивого научного интереса к проблеме.

Отдельную ветвь составляют многомодальные подходы к распознаванию и идентификации,
в которых объединяются комплементарные источники информации: лицо, силуэт,
походка, атрибуты, контекст наблюдения и др.~\cite{Irfan2022multimodal,Koo2018multimodal,Liao2023mmmgcn,Shoukry2022multimodal,Zhao2017spindle}. Предлагаются
как схемы объединения на уровне признаков (конкатенация, линейное объединение),
так и более сложные механизмы на уровне оценок и решающих правил. Вероятностные
и байесовские постановки~\cite{Chen2012bayesface,Liong2014bayesreid,Liu2018bayesreid,Moghaddam2000bayesface,Krivenko2020bayes,Saburov2024bayes} позволяют интерпретировать процесс
принятия решения как объединение нескольких источников свидетельств с учётом
неопределённости, что особенно важно для ИПС, работающих в изменчивых сценах
и при неполной доступности модальностей.

Состояние области показывает, что основные научные и инженерные трудности
связаны с тремя узкими местами. Во‑первых, качество результата существенно
зависит от первичного детектирования объектов, особенно если в контуре ReID
используются как силуэт, так и лицо. Во‑вторых, требуется формировать
устойчивые и одновременно дискриминативные эмбеддинги для каждой модальности,
сохраняя сопоставимость представлений в пределах одной системы принятия решения.
В‑третьих, необходима такая интеграция модальностей, которая была бы
обоснованной, интерпретируемой и корректно работала бы в условиях пропуска
или деградации одной из модальностей.

Следовательно, несмотря на значительное число работ, сохраняется потребность
в методически целостном решении, ориентированном именно на информационно‑поисковые
системы, где результат определяется не только точностью отдельного алгоритма,
но и согласованностью процессов детектирования, формирования эмбеддингов и
интеграции признаков в рамках общего контура принятия решения.

\textbf{Научно‑техническая задача, решаемая в работе.}

Научно‑техническая задача, решаемая в диссертации, состоит в повышении
эффективности информационных процессов автоматического распознавания и
реидентификации людей в информационно‑поисковых системах в условиях
информационной неопределённости и изменяющихся условий наблюдения за счёт
согласованного улучшения детектирования связанных объектов, формирования
устойчивых признаковых представлений и интеграции модальностей при принятии
решения.

\textbf{Цель исследования.}

Целью исследования является повышение эффективности процессов распознавания и
реидентификации людей в информационно‑поисковых системах путём разработки
комплексной методики, обеспечивающей повышение точности идентификации в режиме,
близком к реальному времени, при рациональных вычислительных затратах за счёт
интеграции лицевых и силуэтных признаков.

\textbf{Объект исследования.}

Объектом исследования являются информационные процессы автоматического
распознавания и реидентификации людей в информационно‑поисковых системах
в режиме, близком к реальному времени.

\textbf{Предмет исследования.}

Предметом исследования являются методы, модели и программно‑алгоритмические
средства повышения эффективности распознавания и реидентификации личности
на основе анализа и интеграции лицевых и силуэтных признаков.

\textbf{Задачи исследования.}

Для достижения поставленной цели решены следующие задачи:
\begin{enumerate}
    \item проведён анализ существующих подходов к детектированию, идентификации
    и реидентификации личности и дана оценка их применимости для
    информационно‑поисковых систем, работающих в условиях динамической сцены
    и неполноты данных;

    \item разработан модифицированный критерий обучения нейросетевого детектора
    YOLOv8 для детектирования связанных объектов <<силуэт--лицо>>, включающий
    дополнительный компонент функции потерь $\mathcal{L}_{\text{face}}$ и
    коэффициент $\lambda_{\text{face}}$, а также проведено исследование влияния
    $\lambda_{\text{face}}$ на показатели качества детекции лиц при сохранении
    стабильности детекции силуэта;

    \item разработан метод формирования дискриминативных эмбеддингов для лица
    и силуэта на основе трансформерной архитектуры и функции потерь ArcFace,
    обеспечивающий устойчивость представлений к изменениям условий наблюдения;

    \item разработан метод интеграции эмбеддингов лица и силуэта и принятия
    решения в задаче ReID на основе байесовской модели и MAP‑критерия,
    включая механизм автоматической настройки ключевого гиперпараметра
    комбинирования для адаптации к изменяющимся условиям наблюдения;

    \item реализована программная модель и проведены вычислительные эксперименты
    на репрезентативных наборах данных (CUHK03, Market‑1501, MARS и др.),
    выполнено сравнение с современными подходами и дана оценка эффективности
    предложенных решений.
\end{enumerate}

\textbf{Методология и методы исследования.}

При проведении исследования использованы методы и средства компьютерного зрения
и машинного обучения, включая глубокие нейросетевые архитектуры для
детектирования объектов на изображениях и в видеопотоках, методы обучения
представлений (эмбеддингов) с применением специализированных функций потерь
метрического обучения (ArcFace), а также методы вероятностного моделирования
и принятия решений (байесовский подход и MAP‑критерий) для интеграции
модальностей. Для верификации и сравнения использованы общепринятые метрики
качества (Precision/Recall, mAP, ROC/PR‑кривые, Rank‑k и др.) и вычислительные
эксперименты на открытых наборах данных, а также сравнение с современными
state‑of‑the‑art подходами.

\textbf{Научная новизна работы.}

Научная новизна работы заключается в следующем:
\begin{enumerate}
    \item разработан метод обучения нейросетевого детектора YOLOv8 в задаче
    детектирования связанных объектов <<силуэт--лицо>>, отличающийся
    модификацией функции потерь за счёт введения компонента
    $\mathcal{L}_{\text{face}}$ и коэффициента $\lambda_{\text{face}}$,
    увеличивающего вклад ошибок детекции лиц, локализованных внутри рамки
    силуэта, что повышает полноту и точность формирования ROI лица для
    последующих этапов идентификации и реидентификации;

    \item разработан метод формирования дискриминативных эмбеддингов лица и
    силуэта, ориентированный на совместное использование модальностей в ИПС
    при их неполной доступности и отличающийся двухканальной схемой с
    модально‑специфическими ViT‑кодировщиками, обучаемыми по ArcFace на общем
    множестве идентичностей, а также процедурой стандартизации эмбеддингов
    (включая $L_2$‑нормировку), обеспечивающей их пригодность для последующей
    вероятностной интеграции;

    \item разработан метод интеграции эмбеддингов лица и силуэта и принятия
    решения в задаче ReID, основанный на байесовской модели объединения
    модальностей и многоклассовом MAP‑критерии, позволяющий корректно принимать
    решение при временном отсутствии или низком качестве одной из модальностей;
    метод отличается введением механизма автоматической настройки ключевого
    гиперпараметра комбинирования и обеспечивает адаптивный вклад модальностей
    в зависимости от условий наблюдения.
\end{enumerate}

\textbf{Теоретическая значимость работы.}

Теоретическая значимость работы заключается в развитии формального подхода
к описанию и оценке информационных процессов распознавания и реидентификации
личности в условиях информационной неопределённости за счёт построения
вероятностной модели интеграции признаковых представлений и обоснования
решающего правила на основе MAP‑критерия.

\textbf{Практическая значимость работы.}

Практическая значимость работы состоит в возможности применения разработанного
методического и программного обеспечения в составе информационно‑поисковых
систем распознавания и реидентификации людей в изменяющихся условиях наблюдения
(освещение, ракурс, окклюзии, изменение числа людей в кадре), а также
в системах визуального мониторинга и безопасности, где требуется повышенная
устойчивость к отсутствию или снижению качества отдельных модальностей и
рациональное использование вычислительных ресурсов.

\textbf{Положения, выносимые на защиту.}

На защиту выносятся следующие положения:
\begin{enumerate}
    \item метод обучения нейросетевого детектора YOLOv8 для детектирования
    связанных объектов <<силуэт--лицо>> с модифицированной функцией потерь,
    включающей компонент $\mathcal{L}_{\text{face}}$ и коэффициент
    $\lambda_{\text{face}}$, обеспечивающий улучшение метрик детекции лиц при
    сохранении стабильности детекции класса <<силуэт>>; экспериментально
    установлено, что значения $\lambda_{\text{face}}$ в диапазоне 2,0--2,7
    обеспечивают максимальные показатели по детекции лиц при отсутствии
    деградации по классу <<силуэт>>;

    \item метод формирования дискриминативных эмбеддингов лица и силуэта на базе
    ViT и ArcFace, обеспечивающий получение устойчивых векторных представлений
    и повышение качества сравнения и идентификации объектов в задаче ReID;

    \item метод интеграции эмбеддингов лица и силуэта на основе байесовской
    модели и MAP‑критерия, обеспечивающий повышение качества реидентификации по
    сравнению с раздельным использованием модальностей и линейным объединением;
    по результатам вычислительных экспериментов на наборе CUHK03 получены
    значения Precision~95,65\%, Recall~95,54\%, F1~94,31\%, Accuracy~93,42\%, а
    по метрике mAP достигнуты значения 95,65\% / 98,3\% / 89,2\% для наборов
    CUHK03 / Market‑1501 / MARS соответственно.
\end{enumerate}

\textbf{Соответствие паспорту научной специальности.}

Диссертационная работа соответствует специальности 2.3.8 <<Информатика
и информационные процессы>> и направлениям исследований паспорта специальности:
\begin{itemize}
    \item п.~1 --- в части разработки компьютерных методов описания и оценки
    информационных процессов распознавания и реидентификации, включая
    математическую постановку задачи и вероятностную модель интеграции
    модальностей;

    \item п.~4 --- в части разработки методов и алгоритмов цифровой обработки
    аудиовизуальной информации, включая детектирование объектов, извлечение
    эмбеддингов и построение процедур принятия решения по изображениям
    и видеоданным;

    \item п.~13 --- в части разработки и применения методов распознавания
    образов, нейросетевых технологий и решающих правил, включая
    модифицированный критерий обучения детектора, метод формирования эмбеддингов
    ViT совместно с ArcFace и многоклассовое MAP‑решение.
\end{itemize}

\textbf{Достоверность и обоснованность результатов исследования.}

Достоверность полученных научных результатов подтверждается согласованностью
теоретических выводов с результатами вычислительных экспериментов, корректным
выбором и применением общепринятых метрик качества, воспроизводимостью
экспериментальной процедуры и сравнением с современными решениями на
стандартных наборах данных (CUHK03, Market‑1501, MARS). Обоснованность
результатов обеспечивается непротиворечивостью исходных допущений, использованием
апробированных методов компьютерного зрения, распознавания образов и
вероятностного моделирования, а также достаточной полнотой анализа публикаций
по исследуемой проблеме.

\textbf{Апробация результатов работы.}

Основные положения диссертационной работы докладывались и обсуждались на
плановых научно‑технических конференциях ИПУ РАН (г. Москва, 2020--2022 гг.),
а также на международных и всероссийских конференциях по информатике,
управлению и компьютерному зрению (серии MLSD, ICCT, УБС, МКПУ, ВСПУ и др.).

\textbf{Внедрение результатов работы.}

Результаты диссертационной работы использованы в деятельности
ФГБУ <<НИЦ <<Институт имени Н.\,Е.\,Жуковского>> в работах по автоматическому
распознаванию и реидентификации объектов. Применение результатов подтверждено
актом о внедрении (утверждён 10.06.2024). В акте отмечено, что разработанное
методическое и программное обеспечение позволяет эффективно детектировать
объекты и преобразовывать их в численные характеристики (эмбеддинги), улучшая
точность идентификации, а также обеспечивая повышение эффективности и
сокращение времени принятия решения при использовании систем визуального
наблюдения и безопасности.

\textbf{Личный вклад автора.}

Основные научные результаты получены лично автором. Постановка задачи
исследования осуществлялась совместно с научным руководителем. Автором
выполнены разработка ключевых алгоритмических решений, реализация программных
модулей, постановка и проведение вычислительных экспериментов, анализ и
интерпретация результатов.

\textbf{Публикации.}

Основные результаты диссертационной работы отражены в публикациях автора.
Непосредственно по теме диссертации результаты изложены в трёх самостоятельных
публикациях автора:
\begin{itemize}
    \item Русаков К.\,Д. <<Байесовская модель объединения признаковых
    представлений в задаче реидентификации личности>> // Управление большими
    системами. 2025. Вып.~115. С.~183--202;

    \item Русаков К.\,Д. <<Байесовский алгоритм классификации в задаче
    реидентификации личности>> // Известия Юго‑Западного государственного
    университета. 2025. Т.~29, №~3. С.~92--108;

    \item Rusakov K.\,D. <<Improved Detection of Related Objects on the Example
    of Human Re-Identification Task>> // Proceedings of the 17th International
    Conference Management of Large-Scale System Development (MLSD). IEEE, 2024.
\end{itemize}

\textbf{Объём и структура работы.}

Диссертация состоит из~введения,
\formbytotal{totalchapter}{глав}{ы}{}{},
заключения и
\formbytotal{totalappendix}{приложен}{ия}{ий}{}.
Полный объём диссертации составляет
\formbytotal{TotPages}{страниц}{у}{ы}{}, включая
\formbytotal{totalcount@figure}{рисун}{ок}{ка}{ков} и
\formbytotal{totalcount@table}{таблиц}{у}{ы}{}.
Список литературы содержит
\formbytotal{citenum}{наименован}{ие}{ия}{ий}.

\textbf{Краткое содержание работы.}

Во введении обоснованы актуальность темы исследования, её теоретическая и
практическая значимость, определены объект и предмет исследования,
сформулированы цель и задачи, приведены элементы научной новизны,
сведения об апробации, внедрении и структуре работы.

Первая глава посвящена анализу теоретических основ распознавания и
реидентификации личности, а также современного состояния области. В разделе~1.1
рассматриваются методы и алгоритмы распознавания человека по лицу, включая
классические подходы (PCA/LDA, каскады Виолы--Джонса, HOG и др.) и современные
глубинные архитектуры, в том числе с функциями потерь типа ArcFace и
механизмами устойчивого формирования лицевых представлений в неконтролируемых
условиях съёмки. В разделе~1.2 анализируются методы реидентификации человека,
выходящие за рамки чисто лицевой биометрии: подходы, основанные на силуэтных,
атрибутивных и многомодальных признаках, одно‑ и двухэтапные схемы поиска,
методы переранжирования результатов и частичного сопоставления. В разделе~1.3
рассматриваются коммерческие системы распознавания человека, отечественные и
зарубежные, анализируются их характеристики, области применения и ограничения.
В разделе~1.4 формулируется математическая постановка задачи распознавания и
реидентификации человека в контуре ИПС, вводятся основные обозначения, целевые
функции и эксплуатационные ограничения (ограниченность ресурсов, требования к
метрикам качества и режиму, близкому к реальному времени). В разделе~1.5
подводятся итоги анализа, обобщаются выявленные недостатки однофакторных
решений и фиксируются требования к детектированию, формированию признаков и их
интеграции.

Вторая глава посвящена разработке комплексной методики распознавания и
реидентификации человека в составе информационно‑поисковых систем. В разделе~2.1
рассматривается метод обучения детектора связанных объектов <<силуэт--лицо>>
на базе архитектуры YOLOv8, анализируются особенности сцен массового наблюдения
и формулируется модифицированная функция потерь, учитывающая вложенную задачу
поиска лица внутри силуэта. В разделе~2.2 разрабатывается метод формирования
дискриминативных эмбеддингов лица и силуэта с использованием двухканальной
схемы ViT‑кодировщиков и функции потерь ArcFace, а также рассматриваются
вопросы стандартизации признаков, выбора размерности и построения устойчивых
векторных представлений. В разделе~2.3 формализуется метод интеграции модальностей
на основе байесовской модели, задаются вероятностные модели эмбеддингов,
выводится MAP‑критерий и описывается механизм автоматической настройки
ключевого гиперпараметра комбинирования. В разделе~2.4 подводятся выводы
по главе и обобщаются достоинства предложенного контура
<<детектирование --- формирование эмбеддингов --- байесовская интеграция>>.

Третья глава посвящена экспериментальной оценке эффективности разработанной
комплексной методики. В разделе~3.1 описаны экспериментальная установка,
серверная инфраструктура и используемые наборы данных. В разделе~3.2 приведены
результаты исследования влияния коэффициента $\lambda_{\text{face}}$ на качество
детектирования лиц и силуэтов, включая анализ метрик Precision, Recall,
mAP@0.5 и mAP@0.5:0.95, матрицы ошибок и примеры работы детектора. В разделе~3.3
представлены результаты оценки метода формирования эмбеддингов и метода
интеграции модальностей: t‑SNE‑визуализации признакового пространства, ROC‑
и PR‑кривые, количественные метрики на наборах CUHK03, Market‑1501 и MARS,
а также сравнение с современными методами re‑ID по метрике mAP. В разделе~3.4
сформулированы выводы по главе, обобщающие достигнутый прирост качества и
подтверждающие соблюдение требований к времени обработки и ресурсоёмкости.

В заключении сформулированы основные результаты и выводы диссертационного
исследования, подтверждающие достижение поставленной цели и практическую
значимость разработанной комплексной методики распознавания и реидентификации
людей в информационно‑поисковых системах.